\section{Kriptografi Modern dan Stream Cipher}

\subsection{Karakteristik Kriptografi Modern}

Kriptografi modern menandai pergeseran paradigma fundamental dari era kriptografi klasik, yang dipicu oleh penemuan dan perkembangan pesat teknologi komputer digital sejak awal dekade 1970-an. Pada sistem komputasi digital, seluruh informasi—baik berupa teks, citra, audio, video, hingga instruksi program—direpresentasikan dalam bentuk bit-bit biner ($0$ dan $1$).

Perbedaan mendasar antara kriptografi klasik dan kriptografi modern mencakup beberapa aspek utama:
\begin{enumerate}
  \item \textbf{Unit Pemrosesan Data}: Algoritma klasik bekerja pada taraf karakter (huruf $A$--$Z$ dengan alfabet 26 huruf), sedangkan algoritma modern beroperasi dalam ranah bit (rangkaian bit $0$ dan $1$ atau blok biner).
  \item \textbf{Kekuatan Pengacakan}: Meskipun algoritma modern tetap mengadopsi dua operasi fundamental yang dirumuskan sejak era klasik—yaitu \textit{substitusi} dan \textit{transposisi (permutasi)}—operasinya dirancang jauh lebih kompleks, berlapis, dan matematis agar tahan terhadap analisis frekuensi dan serangan komputasi modern.
  \item \textbf{Operasi Biner Dominan}: Operasi logika \textit{Exclusive-OR} (XOR, dinotasikan $\oplus$) menjadi operasi paling vital di dalam algoritma kriptografi modern karena sifat aljabar dan efisiensi eksekusi perangkat kerasnya.
\end{enumerate}

\subsection{Representasi Data Biner dan Operasi XOR}

Dalam komputasi digital, data disusun dalam hierarki satuan:
\begin{itemize}
  \item \textbf{Bit}: Unit terkecil data digital, bernilai $0$ atau $1$.
  \item \textbf{Byte}: Rangkaian 8 bit. Satu byte dapat merepresentasikan bilangan bulat tak bertanda dari $0$ hingga $255$ ($2^8 - 1$), atau dinyatakan secara ringkas dengan dua digit heksadesimal ($00_{16}$ sampai $\text{FF}_{16}$).
  \item \textbf{Word}: Sekelompok bit berukuran tetap yang diproses prosesor secara atomik (biasanya 32 bit atau 64 bit).
\end{itemize}

\subsubsection{Sifat-Sifat Aljabar Operasi XOR}

Operasi XOR ($\oplus$) bekerja pada dua bit masukan dengan tabel kebenaran: $0 \oplus 0 = 0$, $0 \oplus 1 = 1$, $1 \oplus 0 = 1$, dan $1 \oplus 1 = 0$. Dalam aljabar abstrak, operasi XOR berkorespondensi dengan penjumlahan modulo 2 pada medan Galois $\mathrm{GF}(2)$.

Operasi XOR memiliki sifat-sifat aljabar penting yang menjadikannya fondasi enkripsi simetris:
\begin{enumerate}
  \item \textbf{Komutatif}: $A \oplus B = B \oplus A$
  \item \textbf{Asosiatif}: $(A \oplus B) \oplus C = A \oplus (B \oplus C)$
  \item \textbf{Elemen Identitas}: $A \oplus 0 = A$
  \item \textbf{Involusi (Invers Terhadap Diri Sendiri)}: $A \oplus A = 0$
  \item \textbf{Pembalikan Sempurna (Reversibilitas)}:
  \begin{equation}
    \text{Jika } C = P \oplus K, \quad \text{maka } C \oplus K = (P \oplus K) \oplus K = P \oplus (K \oplus K) = P \oplus 0 = P
  \end{equation}
\end{enumerate}

Sifat inilah yang memungkinkan fungsi enkripsi dan dekripsi menggunakan operasi yang persis sama. Jika plainteks $P$ di-XOR dengan bit kunci $K$, maka hasilnya adalah cipherteks $C$. Untuk mengembalikan plainteks semula, penerima cukup meng-XOR-kan kembali cipherteks $C$ dengan kunci $K$ yang sama.

\begin{figure}[htbp]
  \centering
  \includegraphics[width=0.7\textwidth]{bab-05/xor-tabel-sifat}
  \caption{Tabel kebenaran dan sifat-sifat dasar operasi XOR dalam kriptografi modern.}
  \label{fig:xor-tabel-sifat}
\end{figure}

\subsubsection{Cipher Sederhana dengan Operasi XOR}

Algoritma simetris paling sederhana mengenkripsi deretan byte plainteks $P = p_1 p_2 \dots p_n$ dengan deretan byte kunci $K = k_1 k_2 \dots k_m$. Apabila panjang kunci $m < n$, kunci diulang secara periodik:
\begin{equation}
  c_i = p_i \oplus k_{(i-1 \bmod m) + 1}
\end{equation}

Meskipun sangat efisien dieksekusi, cipher XOR dengan kunci berulang memiliki kelemahan yang serupa dengan Vigen\`ere cipher klasik: rentan terhadap analisis periodisitas (Kasiski) dan analisis frekuensi byte. Oleh karena itu, kriptografi modern mengembangkan algoritma yang membangkitkan deretan kunci yang panjang dan berkarakteristik pseudo-acak (disebut \textit{keystream}).

\begin{lstlisting}[style=pystyle, caption={Implementasi enkripsi dan dekripsi berkas dengan cipher XOR sederhana}]
def xor_cipher(data: bytes, key: bytes) -> bytes:
    """Enkripsi/dekripsi data byte menggunakan kunci byte secara periodik."""
    key_len = len(key)
    return bytes([b ^ key[i % key_len] for i, b in enumerate(data)])

# Contoh pengujian
plaintext = b"Kriptografi Modern ITB 2026"
key = b"KUNCI"
ciphertext = xor_cipher(plaintext, key)
decrypted = xor_cipher(ciphertext, key)

assert decrypted == plaintext
print("Plaintext :", plaintext)
print("Ciphertext (hex):", ciphertext.hex())
print("Decrypted :", decrypted.decode('utf-8'))
\end{lstlisting}

\subsection{Konsep dan Arsitektur Stream Cipher}

Algoritma kriptografi simetris modern untuk data digital diklasifikasikan ke dalam dua keluarga besar:
\begin{enumerate}
  \item \textbf{Stream Cipher (Cipher Alir)}: Memproses plainteks bit demi bit atau byte demi byte secara berurutan pada satu satuan waktu.
  \item \textbf{Block Cipher (Cipher Blok)}: Memproses sekumpulan bit berukuran tetap (misalnya blok 64 bit atau 128 bit) sekaligus secara bersamaan.
\end{enumerate}

\begin{figure}[htbp]
  \centering
  \includegraphics[width=0.85\textwidth]{bab-05/stream-cipher-diagram}
  \caption{Diagram umum stream cipher: pembentukan keystream $k_i$ dari kunci rahasia $K$ dan operasi bitwise XOR.}
  \label{fig:stream-cipher-diagram}
\end{figure}

Pada stream cipher, pembangkit kunci alir (\textit{keystream generator}) menghasilkan deretan bit acak-semu $k_1, k_2, k_3, \dots$ dari sebuah kunci primer rahasia $K$ (dan sering kali ditambah sebuah vektor inisialisasi / \textit{nonce} $IV$). Proses enkripsi dan dekripsi dinyatakan sebagai:
\begin{align}
  c_i &= p_i \oplus k_i \quad \text{(Enkripsi)} \\
  p_i &= c_i \oplus k_i \quad \text{(Dekripsi)}
\end{align}

\begin{table}[htbp]
  \centering
  \small
  \begin{tabular}{lll}
    \toprule
    \textbf{Parameter} & \textbf{Stream Cipher} & \textbf{Block Cipher} \\
    \midrule
    Unit enkripsi & Bit atau byte berurutan & Blok tetap (64, 128 bit) \\
    Kebutuhan memori & Sangat rendah (efisien perangkat keras) & Memerlukan buffer memori per blok \\
    Kecepatan eksekusi & Sangat tinggi, latensi minimal & Bergantung pada kompleksitas round \\
    Kebutuhan padding & Tidak memerlukan padding & Memerlukan padding (PKCS\#7) \\
    Propagasi galat (bit error) & Tidak merambat ke bit lain (synchronous) & Merambat dalam satu blok/antar blok \\
    Domain penerapan utama & Komunikasi streaming, seluler, IoT & Enkripsi penyimpanan berkas, database \\
    \bottomrule
  \end{tabular}
  \caption{Perbandingan mendasar antara Stream Cipher dan Block Cipher.}
  \label{tab:stream-vs-block}
\end{table}

\subsubsection{Synchronous vs Self-Synchronizing Stream Cipher}

Berdasarkan ketergantungan keystream terhadap teks sandi, stream cipher terbagi menjadi dua kategori:
\begin{enumerate}
  \item \textbf{Synchronous Stream Cipher}: Aliran bit kunci $k_i$ dibangkitkan secara independen dari teks asal ($p_i$) maupun teks sandi ($c_i$). Pembangkit keystream hanya bergantung pada kunci primer $K$ dan status internal generator.
  \begin{itemize}
    \item \textit{Keuntungan}: Jika satu bit cipherteks rusak (berubah nilai) selama transmisi, kerusakan tersebut hanya mempengaruhi satu bit plainteks yang bersesuaian saat didekripsi, tanpa merusak bit-bit sesudahnya.
    \item \textit{Kelemahan}: Pengirim dan penerima harus selalu berada dalam sinkronisasi clock yang presisi. Kehilangan satu bit cipherteks akan merusak seluruh dekripsi berikutnya sampai sinkronisasi diatur ulang.
  \end{itemize}

  \item \textbf{Self-Synchronizing (Asynchronous) Stream Cipher}: Setiap bit kunci $k_i$ dibangkitkan berdasarkan fungsi dari sejumlah $n$ bit cipherteks sebelumnya ($c_{i-1}, c_{i-2}, \dots, c_{i-n}$).
  \begin{itemize}
    \item \textit{Keuntungan}: Mampu menyinkronkan kembali dirinya secara otomatis setelah menerima $n$ bit cipherteks yang valid, meskipun beberapa bit transmisi sebelumnya sempat hilang atau terhapus.
    \item \textit{Kelemahan}: Satu bit transmisi yang salah akan merambat dan merusak dekripsi sebanyak $n$ bit berikutnya.
  \end{itemize}
\end{enumerate}

\begin{figure}[htbp]
  \centering
  \begin{minipage}{0.48\textwidth}
    \centering
    \includegraphics[width=\textwidth]{bab-05/synchronous-stream-cipher}
    \caption{Synchronous stream cipher: keystream independen dari cipherteks.}
    \label{fig:sync-stream}
  \end{minipage}
  \hfill
  \begin{minipage}{0.48\textwidth}
    \centering
    \includegraphics[width=\textwidth]{bab-05/self-synchronizing-stream}
    \caption{Self-synchronizing stream cipher: keystream diturunkan dari cipherteks sebelumnya.}
    \label{fig:self-sync-stream}
  \end{minipage}
\end{figure}

\subsection{Linear Feedback Shift Register (LFSR)}

Dalam perancangan perangkat keras, pembangkit keystream paling populer dibangun menggunakan \textbf{Feedback Shift Register (FSR)}. Sebuah shift register terdiri atas deretan sel flip-flop (masing-masing menyimpan $1$ bit) yang digeser serentak ke arah keluaran pada setiap pulsa clock. Bit masukan baru dibentuk melalui fungsi umpan balik (\textit{feedback function}) dari bit-bit isi sel tertentu (disebut \textit{tap}).

Jika fungsi umpan balik hanya menggunakan operasi linear (penjumlahan modulo 2 / XOR), piranti tersebut dinamakan \textbf{Linear Feedback Shift Register (LFSR)}.

\begin{figure}[htbp]
  \centering
  \includegraphics[width=0.75\textwidth]{bab-05/fsr-shift-register}
  \caption{Skema Feedback Shift Register (FSR) dengan tap umpan balik.}
  \label{fig:fsr-shift-register}
\end{figure}

\subsubsection{Formulasi Matematis dan Polinomial Karakteristik}

Sebuah LFSR $n$-tingkat didefinisikan oleh status awal (keadaan \textit{seed}) $s_0, s_1, \dots, s_{n-1}$ dan koefisien umpan balik $c_1, c_2, \dots, c_n \in \{0, 1\}$. Pada setiap detak clock, bit keluaran baru $s_m$ dihasilkan melalui persamaan kekambuhan linear:
\begin{equation}
  s_m = \left( \sum_{i=1}^n c_i s_{m-i} \right) \bmod 2 = (c_1 s_{m-1} \oplus c_2 s_{m-2} \oplus \dots \oplus c_n s_{m-n})
\end{equation}

Perilaku LFSR direpresentasikan melalui \textbf{polinomial karakteristik} atau polinomial umpan balik:
\begin{equation}
  C(x) = 1 + c_1 x + c_2 x^2 + \dots + c_n x^n \in \mathrm{GF}(2)[x]
\end{equation}

\begin{konsep}
\textbf{Teorema Periode Maksimum LFSR}: Sebuah LFSR dengan panjang $n$ sel memiliki periode maksimum (\textit{maximum-length sequence} atau m-sequence) sebesar:
\begin{equation}
  T_{\max} = 2^n - 1
\end{equation}
jika dan hanya jika polinomial karakteristiknya $C(x)$ merupakan \textbf{polinomial primitif} atas medan $\mathrm{GF}(2)$. Keadaan seluruhnya nol ($0,0,\dots,0$) dikecualikan karena akan mengunci keluaran pada deretan nol permanen.
\end{konsep}

\subsubsection{Kelemahan Kriptografis LFSR Murni}

Meskipun deretan bit keluaran LFSR berderajat primitif memenuhi kriteria keacakan statistik Golomb, \textbf{LFSR murni tidak aman untuk digunakan langsung sebagai stream cipher kriptografi}. Karena relasi pergeserannya sepenuhnya bersifat linear, jika penyerang berhasil memperoleh $2n$ bit aliran keluaran berturut-turut melalui serangan \textit{known-plaintext}, seluruh koefisien tap umpan balik $c_1, \dots, c_n$ dapat direkonstruksi secara instan menggunakan \textbf{Algoritma Berlekamp-Massey} dengan kompleksitas waktu polinomial $\mathcal{O}(n^2)$.

Oleh karena itu, dalam cipher alir praktis, LFSR digabungkan secara non-linear (misalnya multi-LFSR dengan clocking tak teratur pada A5/1 atau penyaring non-linear).

\subsection{Algoritma RC4}

\textbf{RC4} (singkatan dari \textit{Rivest Cipher 4} atau secara informal \textit{Ron's Code}) adalah cipher alir berbasis perangkat lunak yang dirancang oleh Prof. Ronald L. Rivest dari Massachusetts Institute of Technology (MIT) pada tahun 1987. Awalnya RC4 dirahasiakan oleh RSA Data Security, Inc., namun spesifikasinya bocor ke publik pada milis Cypherpunks tahun 1994.

RC4 menjadi cipher alir paling populer di dunia pada masanya karena algoritmanya sangat ringkas, cepat, berorientasi byte (sangat cocok untuk CPU 8-bit hingga 64-bit), dan tidak memerlukan manipulasi bit yang rumit. RC4 digunakan secara masif dalam protokol keamanan IEEE 802.11 WEP (\textit{Wired Equivalent Privacy}), WPA, SSL 3.0, TLS 1.0--1.2, Remote Desktop Protocol (RDP), dan BitTorrent.

\begin{figure}[htbp]
  \centering
  \includegraphics[width=0.28\textwidth]{bab-05/potret-ron-rivest}
  \caption{Prof. Ronald L. Rivest, pencipta RC4, RSA, MD5, RC5, dan RC6.}
  \label{fig:ron-rivest}
\end{figure}

\subsubsection{Struktur Algoritma RC4}

RC4 bekerja dengan memelihara sebuah larik status internal (state vector) $S = [S[0], S[1], \dots, S[255]]$ yang merupakan permutasi dari 256 nilai byte unik ($0$ sampai $255$). Kunci rahasia $K$ memiliki panjang fleksibel antara 1 sampai 256 byte (lazimnya 5 sampai 16 byte, yaitu 40 hingga 128 bit).

Algoritma RC4 terdiri atas dua fase:
\begin{enumerate}
  \item \textbf{Key-Scheduling Algorithm (KSA)}: Menginisialisasi dan mengacak larik status $S$ berdasarkan kunci rahasia $K$.
  \item \textbf{Pseudo-Random Generation Algorithm (PRGA)}: Memodifikasi larik status $S$ secara dinamis dan menghasilkan satu byte keystream $z$ pada setiap iterasi.
\end{enumerate}

\begin{figure}[htbp]
  \centering
  \includegraphics[width=0.7\textwidth]{bab-05/rc4-ksa-diagram}
  \caption{Diagram alir Key-Scheduling Algorithm (KSA) pada inisialisasi larik status RC4.}
  \label{fig:rc4-ksa}
\end{figure}

\begin{lstlisting}[style=cstyle, caption={Algoritma KSA dan PRGA pada RC4 dalam Bahasa C}]
void rc4_ksa(uint8_t S[256], const uint8_t *key, int key_len) {
    int i, j = 0;
    // Inisialisasi larik S dengan permutasi identitas
    for (i = 0; i < 256; i++) {
        S[i] = (uint8_t)i;
    }
    // Pengacakan larik S berdasarkan kunci
    for (i = 0; i < 256; i++) {
        j = (j + S[i] + key[i % key_len]) % 256;
        // Tukar nilai S[i] dan S[j]
        uint8_t temp = S[i];
        S[i] = S[j];
        S[j] = temp;
    }
}

uint8_t rc4_prga_byte(uint8_t S[256], int *i_ptr, int *j_ptr) {
    *i_ptr = (*i_ptr + 1) % 256;
    *j_ptr = (*j_ptr + S[*i_ptr]) % 256;

    // Tukar nilai S[i] dan S[j]
    uint8_t temp = S[*i_ptr];
    S[*i_ptr] = S[*j_ptr];
    S[*j_ptr] = temp;

    // Hasilkan byte keystream
    uint8_t t = (S[*i_ptr] + S[*j_ptr]) % 256;
    return S[t];
}
\end{lstlisting}

\begin{figure}[htbp]
  \centering
  \begin{minipage}{0.48\textwidth}
    \centering
    \includegraphics[width=\textwidth]{bab-05/rc4-prga-diagram}
    \caption{Diagram alir PRGA RC4 untuk membangkitkan byte keystream $z$.}
    \label{fig:rc4-prga}
  \end{minipage}
  \hfill
  \begin{minipage}{0.48\textwidth}
    \centering
    \includegraphics[width=\textwidth]{bab-05/rc4-operasi-lengkap}
    \caption{Skema operasional lengkap enkripsi dan dekripsi berbasis RC4.}
    \label{fig:rc4-lengkap}
  \end{minipage}
\end{figure}

\subsubsection{Kelemahan dan Kriptanalisis RC4}

Meskipun sangat populer, penelitian kriptanalisis bertahun-tahun menemukan sejumlah cacat statistik struktural pada RC4:
\begin{enumerate}
  \item \textbf{Invariance Weakness dan Bias Byte Awal}: Scott Fluhrer, Itsik Mantin, dan Adi Shamir (2001) menemukan serangan FMS terhadap RC4. Byte-byte awal keystream memiliki korelasi statistik yang kuat dengan sejumlah kecil byte kunci rahasia. Pada implementasi WEP di Wi-Fi, kunci statis digabungkan dengan IV 24-bit publik ($IV \parallel K$), memungkinkan penyerang pasif merekonstruksi kunci WEP dalam hitungan menit hanya dengan menangkap beberapa ratus ribu paket jaringan.
  \item \textbf{Byte Bias pada TLS}: Peneliti (AlFardan dkk., 2013) mendemonstrasikan serangan Royal Holloway terhadap TLS yang menggunakan RC4, di mana bias pada byte kedua ($z_2$) dan pola byte lainnya dapat mengekstraksi cookie otentikasi sesi HTTPS.
\end{enumerate}

\begin{catatan}
Sejak publikasi RFC 7465 (\textit{Prohibiting RC4 Cipher Suites}) oleh IETF pada tahun 2015, penggunaan RC4 dilarang secara resmi di seluruh versi protokol TLS/SSL.
\end{catatan}

\subsection{Algoritma A5/1 (Enkripsi Seluler GSM)}

\textbf{A5/1} adalah algoritma cipher alir yang digunakan untuk menjamin kerahasiaan komunikasi suara nirkabel pada standar telepon seluler GSM (\textit{Global System for Mobile Communications}) generasi kedua (2G). Algoritma ini dirancang pada tahun 1989 dan awalnya dirahasiakan oleh asosiasi GSM, namun berhasil di-\textit{reverse engineer} dan dipublikasikan oleh Marc Briceno pada 1999.

\begin{figure}[htbp]
  \centering
  \includegraphics[width=0.85\textwidth]{bab-05/a51-gsm-arsitektur}
  \caption{Arsitektur cipher alir A5/1: tiga register LFSR dengan aturan clocking mayoritas.}
  \label{fig:a51-arsitektur}
\end{figure}

\subsubsection{Arsitektur Tiga Register A5/1}

A5/1 menghasilkan aliran bit kunci per percakapan (disebut \textit{burst}, 228 bit per 4.615 ms: 114 bit arah transmisi uplink dan 114 bit arah downlink). Masukan algoritma terdiri dari kunci sesi rahasia 64 bit $K_c$ dan nomor frame publik 22 bit $F_n$.

Algoritma A5/1 memanfaatkan tiga LFSR dengan panjang yang relatif prima:
\begin{enumerate}
  \item \textbf{Register $R_1$}: 19 bit (polinomial $x^{19} + x^5 + x^2 + x + 1$), bit clock pada sel ke-$8$.
  \item \textbf{Register $R_2$}: 22 bit (polinomial $x^{22} + x + 1$), bit clock pada sel ke-$10$.
  \item \textbf{Register $R_3$}: 23 bit (polinomial $x^{23} + x^{15} + x^2 + x + 1$), bit clock pada sel ke-$10$.
\end{enumerate}
Total panjang gabungan ketiga register adalah $19 + 22 + 23 = 64$ bit status.

\subsubsection{Aturan Clocking Mayoritas (Majority Rule)}

Keacakan non-linear A5/1 diperoleh melalui mekanisme clocking tak beraturan:
\begin{equation}
  \text{maj}(x_8, y_{10}, z_{10}) = (x_8 \cdot y_{10}) \oplus (x_8 \cdot z_{10}) \oplus (y_{10} \cdot z_{10})
\end{equation}
Pada setiap detak, nilai bit kontrol ketiga register diperiksa. Register yang bit kontrolnya bernilai sama dengan nilai mayoritas akan digeser (\textit{clocked}), sedangkan register yang bernilai minoritas ditahan (diam). Akibatnya, pada setiap langkah setidaknya dua register (atau ketiga-tiganya) akan bergeser.

Bit keluaran keystream pada setiap detak merupakan hasil XOR bit teratas dari ketiga register:
\begin{equation}
  s_i = R_1[18] \oplus R_2[21] \oplus R_3[22]
\end{equation}

\subsubsection{Status Keamanan A5/1 dan A5/2}

Varian \textbf{A5/2} sengaja dirancang lemah untuk memenuhi regulasi ekspor intelijen pada masanya, dan berhasil dipecahkan secara instan dalam hitungan milidetik. A5/1 sendiri rentan terhadap serangan \textit{time-memory trade-off} menggunakan \textit{rainbow tables} (Barkan, Biham, Shamir 2003; Karsten Nohl 2009). Jaringan 3G dan 4G menggantikan keluarga A5 dengan algoritma berbasis block cipher KASUMI (A5/3) dan SNOW 3G / ZUC.

\subsection{Algoritma Trivium}

\textbf{Trivium} adalah cipher alir berorientasi perangkat keras yang dikembangkan oleh Christophe De Canni\`ere dan Bart Preneel pada tahun 2005 sebagai bagian dari proyek standardisasi Eropa \textbf{eSTREAM} (Profile 2). Trivium dirancang untuk meminimalkan kompleksitas gerbang logika hardware (hanya membutuhkan sekitar 3.000 gerbang NAND) sembari menghasilkan throughput tinggi.

\begin{figure}[htbp]
  \centering
  \includegraphics[width=0.85\textwidth]{bab-05/trivium-arsitektur}
  \caption{Diagram arsitektur cipher alir Trivium dengan 3 register geser non-linear terkopel.}
  \label{fig:trivium-arsitektur}
\end{figure}

Trivium menerima kunci rahasia 80-bit dan IV 80-bit. Status internalnya berukuran total 288 bit, yang dibagi menjadi tiga register geser sirkular tak-linear berukuran 93 bit, 84 bit, dan 111 bit. Umpan balik pada setiap register tidak hanya menggunakan XOR linear, melainkan menyertakan gerbang perkalian biner (AND) dari dua bit status:
\begin{align}
  t_1 &\leftarrow s_{66} \oplus s_{93} \oplus (s_{91} \cdot s_{92}) \oplus s_{171} \\
  t_2 &\leftarrow s_{162} \oplus s_{177} \oplus (s_{175} \cdot s_{176}) \oplus s_{264} \\
  t_3 &\leftarrow s_{243} \oplus s_{288} \oplus (s_{286} \cdot s_{287}) \oplus s_{69}
\end{align}
Bit keystream $z_i$ adalah penjumlahan XOR $z_i = t_1 \oplus t_2 \oplus t_3$. Meskipun strukturnya tampak sangat sederhana, hingga kini belum ada serangan kriptanalisis praktis yang mampu memecahkan Trivium.

\subsection{Salsa20 dan ChaCha20}

Sebagai respons terhadap keusangan RC4 dan kerentanan tabel S-Box terhadap serangan waktu saluran samping (\textit{cache-timing side-channel attacks}), Prof. Daniel J. Bernstein (DJB) merancang keluarga cipher alir berbasis arsitektur \textbf{ARX} (\textit{Addition-Rotation-XOR}): \textbf{Salsa20} pada tahun 2005 (terpilih dalam eSTREAM Profile 1) dan versi penyempurnaannya \textbf{ChaCha20} pada tahun 2008.

\begin{figure}[htbp]
  \centering
  \includegraphics[width=0.3\textwidth]{bab-05/potret-djb-bernstein}
  \caption{Daniel J. Bernstein, kriptografer perancang Salsa20, ChaCha20, Curve25519, dan Ed25519.}
  \label{fig:djb-bernstein}
\end{figure}

\subsubsection{Filosofi Arsitektur ARX Tanpa Tabel Lookup}

Kelemahan paling krusial pada cipher tradisional yang menggunakan tabel substitusi (seperti S-box DES atau AES dalam perangkat lunak) adalah waktu akses memori yang bervariasi bergantung pada indeks yang dituju. Penyerang yang berbagi CPU atau cache dapat menganalisis perbedaan latensi (\textit{timing attack}) untuk merekonstruksi kunci rahasia.

Salsa20 dan ChaCha20 mengeliminasi total tabel pencarian memori. Algoritma ini murni menggunakan tiga operasi dasar prosesor 32-bit:
\begin{enumerate}
  \item Penjumlahan modulo $2^{32}$ (notasi $\boxplus$)
  \item Operasi bitwise XOR (notasi $\oplus$)
  \item Rotasi bit konstan ke kiri (notasi $\lll$)
\end{enumerate}
Ketiga operasi ini selalu dieksekusi dalam waktu konstan (\textit{constant-time}) pada arsitektur prosesor modern mana pun, menjamin kekebalan mutlak dari serangan \textit{cache-timing}.

\subsubsection{Matriks Status dan Fungsi Quarter-Round}

Status internal Salsa20 dan ChaCha20 berukuran 512 bit, yang disusun sebagai matriks $4 \times 4$ dari 16 kata 32-bit:
\begin{equation}
  \mathbf{S} = \begin{pmatrix}
    c_0 & c_1 & c_2 & c_3 \\
    k_0 & k_1 & k_2 & k_3 \\
    k_4 & k_5 & k_6 & k_7 \\
    b_0 & b_1 & n_0 & n_1
  \end{pmatrix}
\end{equation}
di mana $c_0..c_3$ adalah konstanta ASCII bernilai \texttt{"expand 32-byte k"}, $k_0..k_7$ adalah kunci rahasia 256 bit, $b_0..b_1$ adalah counter 64-bit, dan $n_0..n_1$ adalah nonce 64-bit.

\begin{figure}[htbp]
  \centering
  \begin{minipage}{0.48\textwidth}
    \centering
    \includegraphics[width=\textwidth]{bab-05/salsa20-quarter-round}
    \caption{Operasi Quarter-Round pada Salsa20.}
    \label{fig:salsa20-qr}
  \end{minipage}
  \hfill
  \begin{minipage}{0.48\textwidth}
    \centering
    \includegraphics[width=\textwidth]{bab-05/chacha20-quarter-round}
    \caption{Operasi Quarter-Round pada ChaCha20 yang menyempurnakan difusi bit.}
    \label{fig:chacha20-qr}
  \end{minipage}
\end{figure}

Inti komputasi adalah operasi \textbf{Quarter-Round} pada empat kata $(a, b, c, d)$. Pada ChaCha20, fungsi quarter-round dirumuskan sebagai:
\begin{align}
  a &\leftarrow a \boxplus b; \quad d \leftarrow (d \oplus a) \lll 16; \\
  c &\leftarrow c \boxplus d; \quad b \leftarrow (b \oplus c) \lll 12; \\
  a &\leftarrow a \boxplus b; \quad d \leftarrow (d \oplus a) \lll 8; \\
  c &\leftarrow c \boxplus d; \quad b \leftarrow (b \oplus c) \lll 7;
\end{align}

Satu putaran ganda (\textit{double-round}) ChaCha20 mengaplikasikan quarter-round pada empat kolom, dilanjutkan pada empat diagonal. Dalam ChaCha20 standar, proses diulang sebanyak 20 putaran (10 double-rounds). Setelah 20 putaran selesai, status akhir dijumlahkan kembali dengan status awal ($\mathbf{S}_{\text{akhir}} \boxplus \mathbf{S}_{\text{awal}}$) untuk menghasilkan 64 byte keystream.

\subsubsection{Adopsi ChaCha20-Poly1305 di Industri Modern}

ChaCha20 dipadukan dengan kode otentikasi pesan Poly1305 menjadi konstruksi \textbf{AEAD ChaCha20-Poly1305} (distandarisasi dalam RFC 8439). Standar ini menjadi cipher simetris resmi wajib pada:
\begin{itemize}
  \item \textbf{TLS 1.3} (RFC 8446) dan Google Chromium / Android sebagai alternatif cepat berkeamanan tinggi bagi perangkat tanpa instruksi akselerasi AES-NI hardware.
  \item \textbf{WireGuard VPN}: Menggunakan ChaCha20-Poly1305 sebagai algoritma enkripsi utama berkinerja tinggi.
  \item \textbf{OpenSSH}: Menggantikan cipher berbasis CBC dan RC4 lama dengan \texttt{chacha20-poly1305@openssh.com}.
\end{itemize}
ks $P$ di-XOR dengan bit kunci $K$, maka hasilnya adalah cipherteks $C$. Untuk mengembalikan plainteks semula, penerima cukup meng-XOR-kan kembali cipherteks $C$ dengan kunci $K$ yang sama.

\begin{figure}[htbp]
  \centering
  \includegraphics[width=0.7\textwidth]{bab-05/xor-tabel-sifat}
  \caption{Tabel kebenaran dan sifat-sifat dasar operasi XOR dalam kriptografi modern.}
  \label{fig:xor-tabel-sifat}
\end{figure}

\subsubsection{Cipher Sederhana dengan Operasi XOR}

Algoritma simetris paling sederhana mengenkripsi deretan byte plainteks $P = p_1 p_2 \dots p_n$ dengan deretan byte kunci $K = k_1 k_2 \dots k_m$. Apabila panjang kunci $m < n$, kunci diulang secara periodik:
\begin{equation}
  c_i = p_i \oplus k_{(i-1 \bmod m) + 1}
\end{equation}

Meskipun sangat efisien dieksekusi, cipher XOR dengan kunci berulang memiliki kelemahan yang serupa dengan Vigen\`ere cipher klasik: rentan terhadap analisis periodisitas (Kasiski) dan analisis frekuensi byte. Oleh karena itu, kriptografi modern mengembangkan algoritma yang membangkitkan deretan kunci yang panjang dan berkarakteristik pseudo-acak (disebut \textit{keystream}).

\begin{lstlisting}[style=pystyle, caption={Implementasi enkripsi dan dekripsi berkas dengan cipher XOR sederhana}]
def xor_cipher(data: bytes, key: bytes) -> bytes:
    """Enkripsi/dekripsi data byte menggunakan kunci byte secara periodik."""
    key_len = len(key)
    return bytes([b ^ key[i % key_len] for i, b in enumerate(data)])

# Contoh pengujian
plaintext = b"Kriptografi Modern ITB 2026"
key = b"KUNCI"
ciphertext = xor_cipher(plaintext, key)
decrypted = xor_cipher(ciphertext, key)

assert decrypted == plaintext
print("Plaintext :", plaintext)
print("Ciphertext (hex):", ciphertext.hex())
print("Decrypted :", decrypted.decode('utf-8'))
\end{lstlisting}

\subsection{Konsep dan Arsitektur Stream Cipher}

Algoritma kriptografi simetris modern untuk data digital diklasifikasikan ke dalam dua keluarga besar:
\begin{enumerate}
  \item \textbf{Stream Cipher (Cipher Alir)}: Memproses plainteks bit demi bit atau byte demi byte secara berurutan pada satu satuan waktu.
  \item \textbf{Block Cipher (Cipher Blok)}: Memproses sekumpulan bit berukuran tetap (misalnya blok 64 bit atau 128 bit) sekaligus secara bersamaan.
\end{enumerate}

\begin{figure}[htbp]
  \centering
  \includegraphics[width=0.85\textwidth]{bab-05/stream-cipher-diagram}
  \caption{Diagram umum stream cipher: pembentukan keystream $k_i$ dari kunci rahasia $K$ dan operasi bitwise XOR.}
  \label{fig:stream-cipher-diagram}
\end{figure}

Pada stream cipher, pembangkit kunci alir (\textit{keystream generator}) menghasilkan deretan bit acak-semu $k_1, k_2, k_3, \dots$ dari sebuah kunci primer rahasia $K$ (dan sering kali ditambah sebuah vektor inisialisasi / \textit{nonce} $IV$). Proses enkripsi dan dekripsi dinyatakan sebagai:
\begin{align}
  c_i &= p_i \oplus k_i \quad \text{(Enkripsi)} \\
  p_i &= c_i \oplus k_i \quad \text{(Dekripsi)}
\end{align}

\begin{table}[htbp]
  \centering
  \small
  \begin{tabular}{lll}
    \toprule
    \textbf{Parameter} & \textbf{Stream Cipher} & \textbf{Block Cipher} \\
    \midrule
    Unit enkripsi & Bit atau byte berurutan & Blok tetap (64, 128 bit) \\
    Kebutuhan memori & Sangat rendah (efisien perangkat keras) & Memerlukan buffer memori per blok \\
    Kecepatan eksekusi & Sangat tinggi, latensi minimal & Bergantung pada kompleksitas round \\
    Kebutuhan padding & Tidak memerlukan padding & Memerlukan padding (PKCS\#7) \\
    Propagasi galat (bit error) & Tidak merambat ke bit lain (synchronous) & Merambat dalam satu blok/antar blok \\
    Domain penerapan utama & Komunikasi streaming, seluler, IoT & Enkripsi penyimpanan berkas, database \\
    \bottomrule
  \end{tabular}
  \caption{Perbandingan mendasar antara Stream Cipher dan Block Cipher.}
  \label{tab:stream-vs-block}
\end{table}

\subsubsection{Synchronous vs Self-Synchronizing Stream Cipher}

Berdasarkan ketergantungan keystream terhadap teks sandi, stream cipher terbagi menjadi dua kategori:
\begin{enumerate}
  \item \textbf{Synchronous Stream Cipher}: Aliran bit kunci $k_i$ dibangkitkan secara independen dari teks asal ($p_i$) maupun teks sandi ($c_i$). Pembangkit keystream hanya bergantung pada kunci primer $K$ dan status internal generator.
  \begin{itemize}
    \item \textit{Keuntungan}: Jika satu bit cipherteks rusak (berubah nilai) selama transmisi, kerusakan tersebut hanya mempengaruhi satu bit plainteks yang bersesuaian saat didekripsi, tanpa merusak bit-bit sesudahnya.
    \item \textit{Kelemahan}: Pengirim dan penerima harus selalu berada dalam sinkronisasi clock yang presisi. Kehilangan satu bit cipherteks akan merusak seluruh dekripsi berikutnya sampai sinkronisasi diatur ulang.
  \end{itemize}
  
  \item \textbf{Self-Synchronizing (Asynchronous) Stream Cipher}: Setiap bit kunci $k_i$ dibangkitkan berdasarkan fungsi dari sejumlah $n$ bit cipherteks sebelumnya ($c_{i-1}, c_{i-2}, \dots, c_{i-n}$).
  \begin{itemize}
    \item \textit{Keuntungan}: Mampu menyinkronkan kembali dirinya secara otomatis setelah menerima $n$ bit cipherteks yang valid, meskipun beberapa bit transmisi sebelumnya sempat hilang atau terhapus.
    \item \textit{Kelemahan}: Satu bit transmisi yang salah akan merambat dan merusak dekripsi sebanyak $n$ bit berikutnya.
  \end{itemize}
\end{enumerate}

\begin{figure}[htbp]
  \centering
  \begin{minipage}{0.48\textwidth}
    \centering
    \includegraphics[width=\textwidth]{bab-05/synchronous-stream-cipher}
    \caption{Synchronous stream cipher: keystream independen dari cipherteks.}
    \label{fig:sync-stream}
  \end{minipage}
  \hfill
  \begin{minipage}{0.48\textwidth}
    \centering
    \includegraphics[width=\textwidth]{bab-05/self-synchronizing-stream}
    \caption{Self-synchronizing stream cipher: keystream diturunkan dari cipherteks sebelumnya.}
    \label{fig:self-sync-stream}
  \end{minipage}
\end{figure}

\subsection{Linear Feedback Shift Register (LFSR)}

Dalam perancangan perangkat keras, pembangkit keystream paling populer dibangun menggunakan \textbf{Feedback Shift Register (FSR)}. Sebuah shift register terdiri atas deretan sel flip-flop (masing-masing menyimpan $1$ bit) yang digeser serentak ke arah keluaran pada setiap pulsa clock. Bit masukan baru dibentuk melalui fungsi umpan balik (\textit{feedback function}) dari bit-bit isi sel tertentu (disebut \textit{tap}).

Jika fungsi umpan balik hanya menggunakan operasi linear (penjumlahan modulo 2 / XOR), piranti tersebut dinamakan \textbf{Linear Feedback Shift Register (LFSR)}.

\begin{figure}[htbp]
  \centering
  \includegraphics[width=0.75\textwidth]{bab-05/fsr-shift-register}
  \caption{Skema Feedback Shift Register (FSR) dengan tap umpan balik.}
  \label{fig:fsr-shift-register}
\end{figure}

\subsubsection{Formulasi Matematis dan Polinomial Karakteristik}

Sebuah LFSR $n$-tingkat didefinisikan oleh status awal (keadaan \textit{seed}) $s_0, s_1, \dots, s_{n-1}$ dan koefisien umpan balik $c_1, c_2, \dots, c_n \in \{0, 1\}$. Pada setiap detak clock, bit keluaran baru $s_m$ dihasilkan melalui persamaan kekambuhan linear:
\begin{equation}
  s_m = \left( \sum_{i=1}^n c_i s_{m-i} \right) \bmod 2 = (c_1 s_{m-1} \oplus c_2 s_{m-2} \oplus \dots \oplus c_n s_{m-n})
\end{equation}

Perilaku LFSR direpresentasikan melalui \textbf{polinomial karakteristik} atau polinomial umpan balik:
\begin{equation}
  C(x) = 1 + c_1 x + c_2 x^2 + \dots + c_n x^n \in \mathrm{GF}(2)[x]
\end{equation}

\begin{konsep}
\textbf{Teorema Periode Maksimum LFSR}: Sebuah LFSR dengan panjang $n$ sel memiliki periode maksimum (\textit{maximum-length sequence} atau m-sequence) sebesar:
\begin{equation}
  T_{\max} = 2^n - 1
\end{equation}
jika dan hanya jika polinomial karakteristiknya $C(x)$ merupakan \textbf{polinomial primitif} atas medan $\mathrm{GF}(2)$. Keadaan seluruhnya nol ($0,0,\dots,0$) dikecualikan karena akan mengunci keluaran pada deretan nol permanen.
\end{konsep}

\subsubsection{Kelemahan Kriptografis LFSR Murni}

Meskipun deretan bit keluaran LFSR berderajat primitif memenuhi kriteria keacakan statistik Golomb, \textbf{LFSR murni tidak aman untuk digunakan langsung sebagai stream cipher kriptografi}. Karena relasi pergeserannya sepenuhnya bersifat linear, jika penyerang berhasil memperoleh $2n$ bit aliran keluaran berturut-turut melalui serangan \textit{known-plaintext}, seluruh koefisien tap umpan balik $c_1, \dots, c_n$ dapat direkonstruksi secara instan menggunakan \textbf{Algoritma Berlekamp-Massey} dengan kompleksitas waktu polinomial $\mathcal{O}(n^2)$. 

Oleh karena itu, dalam cipher alir praktis, LFSR digabungkan secara non-linear (misalnya multi-LFSR dengan clocking tak teratur pada A5/1 atau penyaring non-linear).

\subsection{Algoritma RC4}

\textbf{RC4} (singkatan dari \textit{Rivest Cipher 4} atau secara informal \textit{Ron's Code}) adalah cipher alir berbasis perangkat lunak yang dirancang oleh Prof. Ronald L. Rivest dari Massachusetts Institute of Technology (MIT) pada tahun 1987. Awalnya RC4 dirahasiakan oleh RSA Data Security, Inc., namun spesifikasinya bocor ke publik pada milis Cypherpunks tahun 1994.

RC4 menjadi cipher alir paling populer di dunia pada masanya karena algoritmanya sangat ringkas, cepat, berorientasi byte (sangat cocok untuk CPU 8-bit hingga 64-bit), dan tidak memerlukan manipulasi bit yang rumit. RC4 digunakan secara masif dalam protokol keamanan IEEE 802.11 WEP (\textit{Wired Equivalent Privacy}), WPA, SSL 3.0, TLS 1.0--1.2, Remote Desktop Protocol (RDP), dan BitTorrent.

\begin{figure}[htbp]
  \centering
  \includegraphics[width=0.28\textwidth]{bab-05/potret-ron-rivest}
  \caption{Prof. Ronald L. Rivest, pencipta RC4, RSA, MD5, RC5, dan RC6.}
  \label{fig:ron-rivest}
\end{figure}

\subsubsection{Struktur Algoritma RC4}

RC4 bekerja dengan memelihara sebuah larik status internal (state vector) $S = [S[0], S[1], \dots, S[255]]$ yang merupakan permutasi dari 256 nilai byte unik ($0$ sampai $255$). Kunci rahasia $K$ memiliki panjang fleksibel antara 1 sampai 256 byte (lazimnya 5 sampai 16 byte, yaitu 40 hingga 128 bit).

Algoritma RC4 terdiri atas dua fase:
\begin{enumerate}
  \item \textbf{Key-Scheduling Algorithm (KSA)}: Menginisialisasi dan mengacak larik status $S$ berdasarkan kunci rahasia $K$.
  \item \textbf{Pseudo-Random Generation Algorithm (PRGA)}: Memodifikasi larik status $S$ secara dinamis dan menghasilkan satu byte keystream $z$ pada setiap iterasi.
\end{enumerate}

\begin{figure}[htbp]
  \centering
  \includegraphics[width=0.7\textwidth]{bab-05/rc4-ksa-diagram}
  \caption{Diagram alir Key-Scheduling Algorithm (KSA) pada inisialisasi larik status RC4.}
  \label{fig:rc4-ksa}
\end{figure}

\begin{lstlisting}[style=cstyle, caption={Algoritma KSA dan PRGA pada RC4 dalam Bahasa C}]
void rc4_ksa(uint8_t S[256], const uint8_t *key, int key_len) {
    int i, j = 0;
    // Inisialisasi larik S dengan permutasi identitas
    for (i = 0; i < 256; i++) {
        S[i] = (uint8_t)i;
    }
    // Pengacakan larik S berdasarkan kunci
    for (i = 0; i < 256; i++) {
        j = (j + S[i] + key[i % key_len]) % 256;
        // Tukar nilai S[i] dan S[j]
        uint8_t temp = S[i];
        S[i] = S[j];
        S[j] = temp;
    }
}

uint8_t rc4_prga_byte(uint8_t S[256], int *i_ptr, int *j_ptr) {
    *i_ptr = (*i_ptr + 1) % 256;
    *j_ptr = (*j_ptr + S[*i_ptr]) % 256;
    
    // Tukar nilai S[i] dan S[j]
    uint8_t temp = S[*i_ptr];
    S[*i_ptr] = S[*j_ptr];
    S[*j_ptr] = temp;
    
    // Hasilkan byte keystream
    uint8_t t = (S[*i_ptr] + S[*j_ptr]) % 256;
    return S[t];
}
\end{lstlisting}

\begin{figure}[htbp]
  \centering
  \begin{minipage}{0.48\textwidth}
    \centering
    \includegraphics[width=\textwidth]{bab-05/rc4-prga-diagram}
    \caption{Diagram alir PRGA RC4 untuk membangkitkan byte keystream $z$.}
    \label{fig:rc4-prga}
  \end{minipage}
  \hfill
  \begin{minipage}{0.48\textwidth}
    \centering
    \includegraphics[width=\textwidth]{bab-05/rc4-operasi-lengkap}
    \caption{Skema operasional lengkap enkripsi dan dekripsi berbasis RC4.}
    \label{fig:rc4-lengkap}
  \end{minipage}
\end{figure}

\subsubsection{Kelemahan dan Kriptanalisis RC4}

Meskipun sangat populer, penelitian kriptanalisis bertahun-tahun menemukan sejumlah cacat statistik struktural pada RC4:
\begin{enumerate}
  \item \textbf{Invariance Weakness dan Bias Byte Awal}: Scott Fluhrer, Itsik Mantin, dan Adi Shamir (2001) menemukan serangan FMS terhadap RC4. Byte-byte awal keystream memiliki korelasi statistik yang kuat dengan sejumlah kecil byte kunci rahasia. Pada implementasi WEP di Wi-Fi, kunci statis digabungkan dengan IV 24-bit publik ($IV \parallel K$), memungkinkan penyerang pasif merekonstruksi kunci WEP dalam hitungan menit hanya dengan menangkap beberapa ratus ribu paket jaringan.
  \item \textbf{Byte Bias pada TLS}: Peneliti (AlFardan dkk., 2013) mendemonstrasikan serangan Royal Holloway terhadap TLS yang menggunakan RC4, di mana bias pada byte kedua ($z_2$) dan pola byte lainnya dapat mengekstraksi cookie otentikasi sesi HTTPS.
\end{enumerate}

\begin{catatan}
Sejak publikasi RFC 7465 (\textit{Prohibiting RC4 Cipher Suites}) oleh IETF pada tahun 2015, penggunaan RC4 dilarang secara resmi di seluruh versi protokol TLS/SSL.
\end{catatan}

\subsection{Algoritma A5/1 (Enkripsi Seluler GSM)}

\textbf{A5/1} adalah algoritma cipher alir yang digunakan untuk menjamin kerahasiaan komunikasi suara nirkabel pada standar telepon seluler GSM (\textit{Global System for Mobile Communications}) generasi kedua (2G). Algoritma ini dirancang pada tahun 1989 dan awalnya dirahasiakan oleh asosiasi GSM, namun berhasil di-\textit{reverse engineer} dan dipublikasikan oleh Marc Briceno pada 1999.

\begin{figure}[htbp]
  \centering
  \includegraphics[width=0.85\textwidth]{bab-05/a51-gsm-arsitektur}
  \caption{Arsitektur cipher alir A5/1: tiga register LFSR dengan aturan clocking mayoritas.}
  \label{fig:a51-arsitektur}
\end{figure}

\subsubsection{Arsitektur Tiga Register A5/1}

A5/1 menghasilkan aliran bit kunci per percakapan (disebut \textit{burst}, 228 bit per 4.615 ms: 114 bit arah transmisi uplink dan 114 bit arah downlink). Masukan algoritma terdiri dari kunci sesi rahasia 64 bit $K_c$ dan nomor frame publik 22 bit $F_n$.

Algoritma A5/1 memanfaatkan tiga LFSR dengan panjang yang relatif prima:
\begin{enumerate}
  \item \textbf{Register $R_1$}: 19 bit (polinomial $x^{19} + x^5 + x^2 + x + 1$), bit clock pada sel ke-$8$.
  \item \textbf{Register $R_2$}: 22 bit (polinomial $x^{22} + x + 1$), bit clock pada sel ke-$10$.
  \item \textbf{Register $R_3$}: 23 bit (polinomial $x^{23} + x^{15} + x^2 + x + 1$), bit clock pada sel ke-$10$.
\end{enumerate}
Total panjang gabungan ketiga register adalah $19 + 22 + 23 = 64$ bit status.

\subsubsection{Aturan Clocking Mayoritas (Majority Rule)}

Keacakan non-linear A5/1 diperoleh melalui mekanisme clocking tak beraturan:
\begin{equation}
  \text{maj}(x_8, y_{10}, z_{10}) = (x_8 \cdot y_{10}) \oplus (x_8 \cdot z_{10}) \oplus (y_{10} \cdot z_{10})
\end{equation}
Pada setiap detak, nilai bit kontrol ketiga register diperiksa. Register yang bit kontrolnya bernilai sama dengan nilai mayoritas akan digeser (\textit{clocked}), sedangkan register yang bernilai minoritas ditahan (diam). Akibatnya, pada setiap langkah setidaknya dua register (atau ketiga-tiganya) akan bergeser.

Bit keluaran keystream pada setiap detak merupakan hasil XOR bit teratas dari ketiga register:
\begin{equation}
  s_i = R_1[18] \oplus R_2[21] \oplus R_3[22]
\end{equation}

\subsubsection{Status Keamanan A5/1 dan A5/2}

Varian \textbf{A5/2} sengaja dirancang lemah untuk memenuhi regulasi ekspor intelijen pada masanya, dan berhasil dipecahkan secara instan dalam hitungan milidetik. A5/1 sendiri rentan terhadap serangan \textit{time-memory trade-off} menggunakan \textit{rainbow tables} (Barkan, Biham, Shamir 2003; Karsten Nohl 2009). Jaringan 3G dan 4G menggantikan keluarga A5 dengan algoritma berbasis block cipher KASUMI (A5/3) dan SNOW 3G / ZUC.

\subsection{Algoritma Trivium}

\textbf{Trivium} adalah cipher alir berorientasi perangkat keras yang dikembangkan oleh Christophe De Canni\`ere dan Bart Preneel pada tahun 2005 sebagai bagian dari proyek standardisasi Eropa \textbf{eSTREAM} (Profile 2). Trivium dirancang untuk meminimalkan kompleksitas gerbang logika hardware (hanya membutuhkan sekitar 3.000 gerbang NAND) sembari menghasilkan throughput tinggi.

\begin{figure}[htbp]
  \centering
  \includegraphics[width=0.85\textwidth]{bab-05/trivium-arsitektur}
  \caption{Diagram arsitektur cipher alir Trivium dengan 3 register geser non-linear terkopel.}
  \label{fig:trivium-arsitektur}
\end{figure}

Trivium menerima kunci rahasia 80-bit dan IV 80-bit. Status internalnya berukuran total 288 bit, yang dibagi menjadi tiga register geser sirkular tak-linear berukuran 93 bit, 84 bit, dan 111 bit. Umpan balik pada setiap register tidak hanya menggunakan XOR linear, melainkan menyertakan gerbang perkalian biner (AND) dari dua bit status:
\begin{align}
  t_1 &\leftarrow s_{66} \oplus s_{93} \oplus (s_{91} \cdot s_{92}) \oplus s_{171} \\
  t_2 &\leftarrow s_{162} \oplus s_{177} \oplus (s_{175} \cdot s_{176}) \oplus s_{264} \\
  t_3 &\leftarrow s_{243} \oplus s_{288} \oplus (s_{286} \cdot s_{287}) \oplus s_{69}
\end{align}
Bit keystream $z_i$ adalah penjumlahan XOR $z_i = t_1 \oplus t_2 \oplus t_3$. Meskipun strukturnya tampak sangat sederhana, hingga kini belum ada serangan kriptanalisis praktis yang mampu memecahkan Trivium.

\subsection{Salsa20 dan ChaCha20}

Sebagai respons terhadap keusangan RC4 dan kerentanan tabel S-Box terhadap serangan waktu saluran samping (\textit{cache-timing side-channel attacks}), Prof. Daniel J. Bernstein (DJB) merancang keluarga cipher alir berbasis arsitektur \textbf{ARX} (\textit{Addition-Rotation-XOR}): \textbf{Salsa20} pada tahun 2005 (terpilih dalam eSTREAM Profile 1) dan versi penyempurnaannya \textbf{ChaCha20} pada tahun 2008.

\begin{figure}[htbp]
  \centering
  \includegraphics[width=0.3\textwidth]{bab-05/potret-djb-bernstein}
  \caption{Daniel J. Bernstein, kriptografer perancang Salsa20, ChaCha20, Curve25519, dan Ed25519.}
  \label{fig:djb-bernstein}
\end{figure}

\subsubsection{Filosofi Arsitektur ARX Tanpa Tabel Lookup}

Kelemahan paling krusial pada cipher tradisional yang menggunakan tabel substitusi (seperti S-box DES atau AES dalam perangkat lunak) adalah waktu akses memori yang bervariasi bergantung pada indeks yang dituju. Penyerang yang berbagi CPU atau cache dapat menganalisis perbedaan latensi (\textit{timing attack}) untuk merekonstruksi kunci rahasia.

Salsa20 dan ChaCha20 mengeliminasi total tabel pencarian memori. Algoritma ini murni menggunakan tiga operasi dasar prosesor 32-bit:
\begin{enumerate}
  \item Penjumlahan modulo $2^{32}$ (notasi $\boxplus$)
  \item Operasi bitwise XOR (notasi $\oplus$)
  \item Rotasi bit konstan ke kiri (notasi $\lll$)
\end{enumerate}
Ketiga operasi ini selalu dieksekusi dalam waktu konstan (\textit{constant-time}) pada arsitektur prosesor modern mana pun, menjamin kekebalan mutlak dari serangan \textit{cache-timing}.

\subsubsection{Matriks Status dan Fungsi Quarter-Round}

Status internal Salsa20 dan ChaCha20 berukuran 512 bit, yang disusun sebagai matriks $4 \times 4$ dari 16 kata 32-bit:
\begin{equation}
  \mathbf{S} = \begin{pmatrix}
    c_0 & c_1 & c_2 & c_3 \\
    k_0 & k_1 & k_2 & k_3 \\
    k_4 & k_5 & k_6 & k_7 \\
    b_0 & b_1 & n_0 & n_1
  \end{pmatrix}
\end{equation}
di mana $c_0..c_3$ adalah konstanta ASCII bernilai \texttt{"expand 32-byte k"}, $k_0..k_7$ adalah kunci rahasia 256 bit, $b_0..b_1$ adalah counter 64-bit, dan $n_0..n_1$ adalah nonce 64-bit.

\begin{figure}[htbp]
  \centering
  \begin{minipage}{0.48\textwidth}
    \centering
    \includegraphics[width=\textwidth]{bab-05/salsa20-quarter-round}
    \caption{Operasi Quarter-Round pada Salsa20.}
    \label{fig:salsa20-qr}
  \end{minipage}
  \hfill
  \begin{minipage}{0.48\textwidth}
    \centering
    \includegraphics[width=\textwidth]{bab-05/chacha20-quarter-round}
    \caption{Operasi Quarter-Round pada ChaCha20 yang menyempurnakan difusi bit.}
    \label{fig:chacha20-qr}
  \end{minipage}
\end{figure}

Inti komputasi adalah operasi \textbf{Quarter-Round} pada empat kata $(a, b, c, d)$. Pada ChaCha20, fungsi quarter-round dirumuskan sebagai:
\begin{align}
  a &\leftarrow a \boxplus b; \quad d \leftarrow (d \oplus a) \lll 16; \\
  c &\leftarrow c \boxplus d; \quad b \leftarrow (b \oplus c) \lll 12; \\
  a &\leftarrow a \boxplus b; \quad d \leftarrow (d \oplus a) \lll 8; \\
  c &\leftarrow c \boxplus d; \quad b \leftarrow (b \oplus c) \lll 7;
\end{align}

Satu putaran ganda (\textit{double-round}) ChaCha20 mengaplikasikan quarter-round pada empat kolom, dilanjutkan pada empat diagonal. Dalam ChaCha20 standar, proses diulang sebanyak 20 putaran (10 double-rounds). Setelah 20 putaran selesai, status akhir dijumlahkan kembali dengan status awal ($\mathbf{S}_{\text{akhir}} \boxplus \mathbf{S}_{\text{awal}}$) untuk menghasilkan 64 byte keystream.

\subsubsection{Adopsi ChaCha20-Poly1305 di Industri Modern}

ChaCha20 dipadukan dengan kode otentikasi pesan Poly1305 menjadi konstruksi \textbf{AEAD ChaCha20-Poly1305} (distandarisasi dalam RFC 8439). Standar ini menjadi cipher simetris resmi wajib pada:
\begin{itemize}
  \item \textbf{TLS 1.3} (RFC 8446) dan Google Chromium / Android sebagai alternatif cepat berkeamanan tinggi bagi perangkat tanpa instruksi akselerasi AES-NI hardware.
  \item \textbf{WireGuard VPN}: Menggunakan ChaCha20-Poly1305 sebagai algoritma enkripsi utama berkinerja tinggi.
  \item \textbf{OpenSSH}: Menggantikan cipher berbasis CBC dan RC4 lama dengan \texttt{chacha20-poly1305@openssh.com}.
\end{itemize}
